% ═══════════════════════════════════════════════════════════════
% MUSTERLÖSUNG: Reaktionsgleichungen aufstellen und ausgleichen
% Thema 2.5 | Chemie Klasse 8 | 2 Seiten | 32 BE
% ═══════════════════════════════════════════════════════════════

\documentclass[10pt, a4paper]{article}

\usepackage[utf8]{inputenc}
\usepackage[T1]{fontenc}
\usepackage[ngerman]{babel}

\usepackage{helvet}
\renewcommand{\familydefault}{\sfdefault}

\usepackage{microtype}
\usepackage[margin=1.3cm, top=1.8cm, bottom=1.2cm]{geometry}
\usepackage{enumitem}
\usepackage{tabularx}
\usepackage{booktabs}
\usepackage{array}
\usepackage{multicol}
\usepackage{tikz}
\usetikzlibrary{calc, shapes.geometric, arrows.meta}

\usepackage{amsmath, amssymb}
\usepackage{siunitx}
\sisetup{locale = DE, per-mode = symbol}

\usepackage{fancyhdr}
\usepackage{lastpage}

\usepackage{xcolor}
\usepackage{tcolorbox}
\tcbuselibrary{skins}

% ═══════════════════════════════════════════════════════════════
% FARBEN
% ═══════════════════════════════════════════════════════════════

\definecolor{chemie}{HTML}{2a7a4b}
\definecolor{hellgrau}{HTML}{F5F5F5}
\definecolor{mittelgrau}{HTML}{888888}
\definecolor{dunkelgrau}{HTML}{444444}
\definecolor{loesung}{HTML}{c62828}
\definecolor{metall}{HTML}{2c5aa0}

\colorlet{fachfarbe}{chemie}

% ═══════════════════════════════════════════════════════════════
% VARIABLEN
% ═══════════════════════════════════════════════════════════════

\newcommand{\fach}{Chemie}
\newcommand{\thema}{Reaktionsgleichungen aufstellen und ausgleichen}
\newcommand{\klasse}{8}
\newcommand{\maxpunkte}{32}

% ═══════════════════════════════════════════════════════════════
% KOPF- UND FUSSZEILE
% ═══════════════════════════════════════════════════════════════

\pagestyle{fancy}
\fancyhf{}
\renewcommand{\headrulewidth}{0.4pt}
\renewcommand{\footrulewidth}{0pt}
\lhead{\small \fach{} Klasse \klasse{} | \thema}
\rhead{\small \textcolor{loesung}{\textbf{MUSTERLÖSUNG}}}
\cfoot{\small\textcolor{mittelgrau}{Seite \thepage{} von \pageref{LastPage}}}

% ═══════════════════════════════════════════════════════════════
% BOXEN
% ═══════════════════════════════════════════════════════════════

\newtcolorbox{merkkasten}{
    colback=hellgrau,
    colframe=hellgrau,
    boxrule=0pt,
    arc=0pt,
    left=6pt, right=6pt, top=5pt, bottom=4pt,
    borderline north={2pt}{0pt}{fachfarbe},
    before skip=4pt, after skip=4pt
}

\newtcolorbox{algorithmusbox}[1][]{
    enhanced,
    colback=fachfarbe!8,
    colframe=fachfarbe,
    boxrule=1.5pt,
    arc=0pt,
    left=8pt, right=8pt, top=12pt, bottom=6pt,
    before skip=10pt, after skip=6pt,
    title={#1},
    fonttitle=\bfseries\small,
    coltitle=white,
    attach boxed title to top left={yshift=-3mm, xshift=4mm},
    boxed title style={colback=fachfarbe, arc=0pt, boxrule=0pt, left=4pt, right=4pt}
}

\newtcolorbox{hinweisbox}{
    colback=metall!10,
    colframe=metall,
    boxrule=1pt,
    arc=0pt,
    left=6pt, right=6pt, top=4pt, bottom=4pt,
    before skip=4pt, after skip=4pt
}

% ═══════════════════════════════════════════════════════════════
% BEFEHLE
% ═══════════════════════════════════════════════════════════════

\newcommand{\lsg}[1]{\textcolor{loesung}{\textbf{#1}}}
\newcommand{\punktebox}[1]{\hfill\small\textcolor{mittelgrau}{[\_\_/#1]}}

\newcommand{\aufgabe}[2]{%
    \vspace{6pt}%
    \noindent\textbf{\textcolor{fachfarbe}{Aufgabe #1}} #2%
    \vspace{3pt}%
    \nopagebreak[4]%
}

\newcommand{\operator}[1]{\textbf{\textcolor{fachfarbe}{#1}}}

\newlist{teil}{enumerate}{2}
\setlist[teil,1]{label=\alph*), leftmargin=1.6em, itemsep=2pt, parsep=0pt, topsep=2pt}

\setlength{\parindent}{0pt}
\setlength{\parskip}{2pt}

% ═══════════════════════════════════════════════════════════════
% DOKUMENT
% ═══════════════════════════════════════════════════════════════

\begin{document}

% ═══════════════════════════════════════════════════════════════
% HEADER
% ═══════════════════════════════════════════════════════════════

\begin{center}
    {\large\bfseries\textcolor{fachfarbe}{\thema}}\\[0.1em]
    {\small Musterlösung | \fach{} Klasse \klasse{} | Thema 2.5}
\end{center}
\vspace{-0.4em}
\hrule
\vspace{0.4em}

% ═══════════════════════════════════════════════════════════════
% WERTIGKEIT + TABELLE
% ═══════════════════════════════════════════════════════════════

\begin{minipage}[t]{0.58\textwidth}
\begin{merkkasten}
\textbf{Wertigkeit – Die „Händezahl"}\\[0.2em]
Die \lsg{Wertigkeit} gibt an, wie viele \lsg{Bindungen} ein Atom eingehen kann.\\[0.3em]
Man kann sie sich als „\lsg{Händezahl}" vorstellen.
\end{merkkasten}
\end{minipage}%
\hfill
\begin{minipage}[t]{0.38\textwidth}
\centering
\small
\textbf{Wertigkeitstabelle}\\[0.3em]
\begin{tabular}{|c|c|c|c|c|c|c|c|}
\hline
\textbf{Na} & \textbf{Mg} & \textbf{Al} & \textbf{Fe} & \textbf{Cu} & \textbf{Zn} & \textbf{Ca} & \textbf{O} \\
\hline
I & II & III & II/III & I/II & II & II & II \\
\hline
\end{tabular}
\end{minipage}

\vspace{4pt}

\begin{hinweisbox}
\small
\textbf{Warum bilden sich diese Formeln?} Bei Metalloxiden gibt das Metall Elektronen ab, Sauerstoff nimmt sie auf. So entstehen \textbf{Ionen}, die sich anziehen (Ionenbindung). Die Wertigkeit zeigt, wie viele Elektronen abgegeben bzw. aufgenommen werden.
\end{hinweisbox}

\vspace{2pt}

% ═══════════════════════════════════════════════════════════════
% ALGORITHMUS 1: VERHÄLTNISFORMEL
% ═══════════════════════════════════════════════════════════════

\begin{algorithmusbox}[Algorithmus 1: Verhältnisformel aufstellen]
\small
\begin{tabular}{@{}c p{12.5cm}@{}}
\textbf{1} & \textbf{Elementsymbole aufschreiben:} \hfill Li \quad O \\[0.4em]
\textbf{2} & \textbf{Wertigkeiten ablesen:} \hfill I \quad II \\[0.4em]
\textbf{3} & \textbf{Kreuzregel anwenden:} Wertigkeit von A → Index von B (und umgekehrt)\\
& \hfill Li\textsubscript{\lsg{2}}O\textsubscript{\lsg{1}} → \textbf{Li\textsubscript{2}O} \\
\end{tabular}
\end{algorithmusbox}

\vspace{-2pt}

% ═══════════════════════════════════════════════════════════════
% ALGORITHMUS 2: REAKTIONSGLEICHUNG
% ═══════════════════════════════════════════════════════════════

\begin{algorithmusbox}[Algorithmus 2: Reaktionsgleichung aufstellen und ausgleichen]
\small
\begin{tabular}{@{}c p{12.5cm}@{}}
\textbf{1} & \textbf{Wortgleichung:} Lithium + Sauerstoff → Lithiumoxid \\[0.4em]
\textbf{2} & \textbf{Formelgleichung:} Li + O\textsubscript{2} → Li\textsubscript{2}O \quad {\scriptsize(Achtung: Sauerstoff ist O\textsubscript{2}!)} \\[0.4em]
\textbf{3} & \textbf{Atome zählen:} Links: Li = \lsg{1}, O = \lsg{2} \quad Rechts: Li = \lsg{2}, O = \lsg{1} \\[0.4em]
\textbf{4} & \textbf{Vorfaktoren ausgleichen:} \\
& \quad a) Produkt verdoppeln, bis O-Zahl gerade: Li\textsubscript{2}O hat 1 O $\rightarrow$ 2 Li\textsubscript{2}O hat \lsg{2} O \\
& \quad b) O\textsubscript{2}-Anzahl = O rechts $\div$ 2 = \lsg{1} \\
& \quad c) Metall-Anzahl = Metall-Atome rechts = \lsg{4} \\
& \quad $\Rightarrow$ \lsg{4} Li + \lsg{1} O\textsubscript{2} $\rightarrow$ \lsg{2} Li\textsubscript{2}O \\
& \quad d) \textbf{Kontrolle:} Li: \lsg{4} = \lsg{4} \checkmark \quad O: \lsg{2} = \lsg{2} \checkmark \\
\end{tabular}
\end{algorithmusbox}

% ═══════════════════════════════════════════════════════════════
% AUFGABE 1
% ═══════════════════════════════════════════════════════════════

\aufgabe{1 ($\star$)}{Wertigkeiten ablesen} \punktebox{4}

\small
\begin{tabular}{|p{3cm}|p{2cm}|p{8cm}|}
\hline
\textbf{Element} & \textbf{Wertigkeit} & \textbf{Bedeutung} \\
\hline
Magnesium (Mg) & \lsg{II} & Mg gibt \lsg{2} Elektron(en) ab. \\
\hline
Aluminium (Al) & \lsg{III} & Al gibt \lsg{3} Elektron(en) ab. \\
\hline
Sauerstoff (O) & \lsg{II} & O nimmt \lsg{2} Elektron(en) auf. \\
\hline
Calcium (Ca) & \lsg{II} & Ca gibt \lsg{2} Elektron(en) ab. \\
\hline
\end{tabular}

\vspace{4pt}

% ═══════════════════════════════════════════════════════════════
% AUFGABE 2
% ═══════════════════════════════════════════════════════════════

\aufgabe{2 ($\star\star$)}{Verhältnisformeln mit der Kreuzregel} \punktebox{6}

\small
\begin{tabular}{|p{4cm}|p{3cm}|p{5cm}|}
\hline
\textbf{Metall + Sauerstoff} & \textbf{Wertigkeiten} & \textbf{Oxidformel} \\
\hline
Magnesium + Sauerstoff & Mg(II), O(II) & MgO \\
\hline
Calcium + Sauerstoff & Ca(\lsg{II}), O(\lsg{II}) & \lsg{CaO} \\
\hline
Zink + Sauerstoff & Zn(\lsg{II}), O(\lsg{II}) & \lsg{ZnO} \\
\hline
Aluminium + Sauerstoff & Al(\lsg{III}), O(\lsg{II}) & \lsg{Al\textsubscript{2}O\textsubscript{3}} \\
\hline
Eisen + Sauerstoff & Fe(III), O(II) & \lsg{Fe\textsubscript{2}O\textsubscript{3}} \\
\hline
Kupfer + Sauerstoff & Cu(II), O(II) & \lsg{CuO} \\
\hline
\end{tabular}

\newpage

% ═══════════════════════════════════════════════════════════════
% AUFGABE 3
% ═══════════════════════════════════════════════════════════════

\aufgabe{3 ($\star\star$)}{Reaktionsgleichungen aufstellen und ausgleichen} \punktebox{12}

\small

\textbf{a) Magnesium} (Wertigkeit: II)

\begin{tabular}{@{}p{7.5cm}p{5.5cm}@{}}
\textbf{Alg. 1:} Mg(\lsg{II}) + O(\lsg{II}) → Oxidformel: \lsg{MgO} & {\scriptsize Gleiche Wertigkeit → 1:1} \\[0.6em]
\textbf{Alg. 2:} Mg + O\textsubscript{2} → \lsg{MgO} & Atome zählen: \\
& Links: Mg = \lsg{1}, O = \lsg{2} \\
& Rechts: Mg = \lsg{1}, O = \lsg{1} \\
Ausgeglichen: \lsg{2} Mg + \lsg{1} O\textsubscript{2} → \lsg{2} \lsg{MgO} & \\
\end{tabular}

\vspace{6pt}

\textbf{b) Aluminium} (Wertigkeit: III)

\begin{tabular}{@{}p{7.5cm}p{5.5cm}@{}}
\textbf{Alg. 1:} Al(\lsg{III}) + O(\lsg{II}) → Oxidformel: \lsg{Al\textsubscript{2}O\textsubscript{3}} & {\scriptsize Kreuzregel: 3→O, 2→Al} \\[0.6em]
\textbf{Alg. 2:} Al + O\textsubscript{2} → \lsg{Al\textsubscript{2}O\textsubscript{3}} & Atome zählen: \\
& Links: Al = \lsg{1}, O = \lsg{2} \\
& Rechts: Al = \lsg{2}, O = \lsg{3} \\
Ausgeglichen: \lsg{4} Al + \lsg{3} O\textsubscript{2} → \lsg{2} \lsg{Al\textsubscript{2}O\textsubscript{3}} & \\
\end{tabular}

\vspace{6pt}

\textbf{c) Eisen} (Wertigkeit: III)

\begin{tabular}{@{}p{7.5cm}p{5.5cm}@{}}
\textbf{Alg. 1:} Fe(\lsg{III}) + O(\lsg{II}) → Oxidformel: \lsg{Fe\textsubscript{2}O\textsubscript{3}} & {\scriptsize Kreuzregel: 3→O, 2→Fe} \\[0.6em]
\textbf{Alg. 2:} Fe + O\textsubscript{2} → \lsg{Fe\textsubscript{2}O\textsubscript{3}} & Atome zählen: \\
& Links: Fe = \lsg{1}, O = \lsg{2} \\
& Rechts: Fe = \lsg{2}, O = \lsg{3} \\
Ausgeglichen: \lsg{4} Fe + \lsg{3} O\textsubscript{2} → \lsg{2} \lsg{Fe\textsubscript{2}O\textsubscript{3}} & \\
\end{tabular}

\vspace{6pt}

\textbf{d) Kupfer} (Wertigkeit: II)

\begin{tabular}{@{}p{7.5cm}p{5.5cm}@{}}
Oxidformel: \lsg{CuO} & {\scriptsize Cu(II), O(II) → 1:1} \\
Ausgeglichen: \lsg{2} Cu + \lsg{1} O\textsubscript{2} → \lsg{2} \lsg{CuO} & \\
\end{tabular}

\vspace{8pt}

% ═══════════════════════════════════════════════════════════════
% AUFGABE 4
% ═══════════════════════════════════════════════════════════════

\aufgabe{4 ($\star\star\star$)}{Fehler finden und korrigieren} \punktebox{6}

\textbf{a) Max schreibt:} \quad 2 Al + 3 O → Al\textsubscript{2}O\textsubscript{3}

\vspace{2pt}
Fehler: \lsg{Sauerstoff muss als O\textsubscript{2} geschrieben werden, nicht als O.}

Korrektur: \lsg{4 Al + 3 O\textsubscript{2} → 2 Al\textsubscript{2}O\textsubscript{3}}

\vspace{4pt}

\textbf{b) Max schreibt:} \quad Mg + O\textsubscript{2} → MgO\textsubscript{2}

\vspace{2pt}
Fehler: \lsg{Max hat den Index geändert statt einen Vorfaktor zu nutzen. MgO\textsubscript{2} ist ein anderer Stoff als MgO!}

Korrektur: \lsg{2 Mg + O\textsubscript{2} → 2 MgO}

\vspace{8pt}

% ═══════════════════════════════════════════════════════════════
% AUFGABE 5
% ═══════════════════════════════════════════════════════════════

\aufgabe{5 ($\star\star\star$)}{Transfer: Unbekanntes Metall} \punktebox{4}

\vspace{4pt}
\begin{tabular}{@{}ll@{}}
Wortgleichung: & Natrium + Sauerstoff → Natriumoxid \\[0.5em]
Oxidformel (Alg. 1): & Na(\lsg{I}) + O(\lsg{II}) → Na\textsubscript{\lsg{2}}O\textsubscript{\lsg{1}} = \lsg{Na\textsubscript{2}O} \\[0.5em]
Ausgeglichene Gleichung (Alg. 2): & \lsg{4} Na + \lsg{1} O\textsubscript{2} → \lsg{2} Na\textsubscript{2}O \\
\end{tabular}

% ═══════════════════════════════════════════════════════════════
% PUNKTE
% ═══════════════════════════════════════════════════════════════

\vfill
\hrule
\vspace{0.3em}
\hfill\textbf{Punkte:} \textcolor{loesung}{32} / \maxpunkte

\end{document}
